\documentclass[10pt,a4paper]{amsart}
\usepackage[margin=19mm]{geometry}
\usepackage{amsmath,amssymb,mathtools}
\usepackage[scr=rsfs,cal=euler]{mathalfa}
\usepackage{xfrac}
\usepackage[unicode,hidelinks,bookmarks=false]{hyperref}
\newcommand{\ie}{i.e.\ }
\newcommand{\cf}{cf.\ }
\newcommand{\resp}{respectively\ }
\newcommand{\Subsection}{\subsection}
\providecommand{\nobreakdash}{\nobreak\hbox{-}}
\setlength{\emergencystretch}{2em}
\multlinegap0pt
\usepackage[shortlabels]{enumitem}
\usepackage{amssymb}
\usepackage{xcolor}
\newtheorem{lem}{Lemma}
\newtheorem{prop}{Proposition}
\def\Aut{\ensuremath{\operatorname{Aut}}}
\def\GL{\ensuremath{\operatorname{GL}}}
\def\Mu{\ensuremath{\mathcal{M}}}
\title{Representability of the automorphism group of finitely generated vertex algebras}
\author{Terry Gannon, Robin Mader, Arturo Pianzola}
\date{Source: arXiv:2605.15605v1 (CC BY 4.0)}
\begin{document}
\maketitle

\noindent\emph{Source and licence.} Representability of the automorphism group of finitely generated vertex algebras. Version arXiv:2605.15605v1 (CC BY 4.0), distributed by arXiv under the Creative Commons Attribution 4.0 International licence, http://creativecommons.org/licenses/by/4.0/, as recorded in the arXiv licence metadata for this exact version and shown on the abstract page https://arxiv.org/abs/2605.15605v1. Full licence notice: \texttt{licenses/arxiv-2605.15605-CC-BY-4.0.txt}.\par\smallskip

\noindent\textit{Editorial scope.} Counted unit: the article's prop prop2 (source0.tex lines 149--160). The article's complete original proof of it is delivered with it (source0.tex lines 162--173, one \texttt{proof} environment of 12 lines). No mathematics is written, repaired or re-derived: the statement and the proof are the article's own lines, in the article's own notation, and no retained line is substituted or re-flowed.  Retained uncounted as the source-local context the proof needs: lines 135--137.  Nothing was omitted from any retained span, so the package carries no omission record.  No retained line contains a citation, so the wrapper carries no bibliography and no bibliography entry.  No retained line contains a figure environment, an image-inclusion command or drawing code, so no image is delivered and the package declares no asset dependency.  Editorial formatting only, in addition: an AMS \texttt{amsart} preamble that restates the article's own declarations and macros used by the retained lines, namely the environments \texttt{lem}, \texttt{prop} on separate counters, together with the standard packages those lines need.  The retained environments print in this document as Lemma 1, Proposition 1 in order of appearance, whereas the article prints the same items with the numbering of its own full document.  The complete native source of the article as distributed in the arXiv e-print for this exact version is bundled unchanged as \texttt{sources/200666/source0.tex}.
\begin{lem}\label{Lemma1}
		The above map $ \theta \mapsto \widetilde \theta$ is an injective homomorphism of groups $ \GL_R(F) \to \Aut (A(X)).$ In fact, this map induces a group isomorphism between $\GL_R(F)$ and  \[ \Aut(A(X),F) = \{ \theta \in \Aut(A(X)) \mid \theta(F) = F\}.\] 	
\end{lem}

\begin{prop}\label{prop2} Let $A$ be an $\Mu$-algebra over $R$ which contains a copy of $F$ as an $R$-submodule. Let $\Aut(A,F) = \{ \theta \in \Aut(A) \mid \theta(F) = F\}.$ If  $F$ generates $A,$ then:

\begin{enumerate}[label = (\roman*)]
	\item \label{p2i} The map 
\begin{equation}\label{tildemap} \Aut(A,F) \to \Aut(A(X), F), \quad \theta \mapsto \widetilde \theta = \widetilde{\theta|_F}. \end{equation}
is a group monomorphism.

\item \label{p2ii}  There exists a unique surjective $\Mu$-algebra homomorphism $\eta\colon A(X) \to A$ that is the identity on $X.$ 

\item \label{p2iii} Let $I$ denote the kernel of $\eta$. The image of the map in \eqref{tildemap} consists of those $\varphi \in \Aut(A(X),F)$ with $\varphi(I) = I$.
\end{enumerate}
\end{prop}

\begin{proof}
	\ref{p2i} The map is clearly a group homomorphism. Since $F$ generates $A$ we can appeal to Lemma \ref{Lemma1} to conclude that it is injective.

	\ref{p2ii} By assumption $A$ has a copy of $F = \bigoplus_{x \in X} Rx.$ We can thus view the elements of $X$ as elements of $A$. Now (ii) is clear by the universal property of $A(X)$ and the assumption that $F$ generates $A.$

	\ref{p2iii} For $\theta \in \Aut(A,F)$ we have $\eta \circ \widetilde \theta = \theta \circ \eta$, because both sides of the equation are $\Mu$-algebra homomorphisms with the same restriction to $X.$
	Thus $\eta \widetilde \theta (I) = \theta \eta (I) = 0$, or equivalently, $\widetilde \theta(I) \subset I$. The same reasoning applies to $\theta^{-1}$ and so $\widetilde \theta(I) = I$.
	Conversely, let $\varphi \in \Aut(A(X),F)$ with $\varphi(I) = I$. 
	Since $\eta$ is surjective,  
	$\varphi$ induces an automorphism $\theta$ on the quotient $A \simeq A(X)/I$ via $\theta \eta = \eta \varphi$.  Since $\eta(F) = F$, we get $\theta(F) = \eta (\varphi(F)) = F$, so that $\theta \in \Aut(A,F)$.
	From $\eta \varphi = \eta \widetilde \theta$ we deduce $\varphi|_F = \widetilde \theta|_F$, whence $\varphi = \widetilde \theta$.
\end{proof}

\end{document}
